% =============================================================================
%  s01_seccion.tex — Sección 1: Introducción / Planteamiento del problema
% =============================================================================

\section{Introducción}
\label{sec:introduccion}

% ───────────────────────────────────────────────────────────────────────────────────────────
\begin{frame}[c]{Pregunta de Investigación}
  \hspace{0.8cm}%
  \begin{minipage}{0.88\textwidth}
    {\color{ColorAcento}\rule{0.5\textwidth}{1.2pt}}\\[0.4cm]
    {\large\bfseries\color{ColorPrincipal}
      Pregunta de Investigación}\\[0.6cm]
    {\normalsize\color{black}
      ¿En qué medida la optimización topológica de la red de transporte 
      de la Ciudad deMéxico y la Zona Metropolitana del Valle de México, 
      evaluada mediante el motor analítico \textit{VFT} \cite{garibaldi2023vft}, 
      demuestra que la implementación de corredores anillares reduce 
      la vulnerabilidad funcional y mejora la eficiencia global del sistema en 
      comparación con el modelo radial vigente?.}\\[0.5cm]
    {\color{ColorAcento}\rule{0.3\textwidth}{0.8pt}}
    \vfill
  \end{minipage}
\end{frame}

% ───────────────────────────────────────────────────────────────────────────────────────────
\begin{frame}[c]{Paráfrasis de la Hipótesis}
  \hspace{0.8cm}%
  \begin{minipage}{0.88\textwidth}
    {\color{ColorAcento}\rule{0.5\textwidth}{1.2pt}}\\[0.4cm]
    {\large\bfseries\color{ColorPrincipal}
        ¿Cuánto mejoraría la movilidad de las personas si el
        Periférico tuviera transporte masivo circular,
        en lugar de
        que todos los viajes tengan que pasar por el centro?}\\[0.5cm]
    {\normalsize\color{black}
      El Anillo Periférico fue construido para conectar la ciudad,
      pero terminó siendo una frontera. Hoy es una autopista
      exclusiva del automóvil privado, mientras las alcaldías
      del sur (Xochimilco, Tláhuac, Tlalpan, Milpa Alta)
      quedan desconectadas del sistema masivo.}\\[0.4cm]
    {\color{ColorAcento}\rule{0.3\textwidth}{0.8pt}}
  \end{minipage}
\end{frame}

% ── Frame: Metodología en 4 fases ────────────────────────────────────────────
\begin{frame}{Metodología en 4 Fases}
  {\small\color{ColorPrincipal}\textit{%
    Enfoque cuantitativo — datos GTFS de la \zmvm\ procesados mediante teoría de grafos
    y redes complejas.}}
  \vspace{0.4cm}
  \begin{center}
  \begin{tikzpicture}[
    fase/.style  = {rectangle, rounded corners=5pt,
                    draw=ColorPrincipal, line width=1pt,
                    fill=ColorPrincipal, text=white,
                    text width=2.6cm, align=center,
                    minimum height=0.85cm, font=\small\bfseries},
    desc/.style  = {text width=2.6cm, align=justify, font=\scriptsize, text=ColorTexto},
    flecha/.style= {-Stealth, thick, color=ColorAcento},
  ]
    \node[fase] (f1) {\textbf{01}\\[1pt]Construir};
    \node[fase, right=0.35cm of f1] (f2) {\textbf{02}\\[1pt]Diagnosticar};
    \node[fase, right=0.35cm of f2] (f3) {\textbf{03}\\[1pt]Simular};
    \node[fase, right=0.35cm of f3] (f4) {\textbf{04}\\[1pt]Proponer};

    \draw[flecha] (f1.east) -- (f2.west);
    \draw[flecha] (f2.east) -- (f3.west);
    \draw[flecha] (f3.east) -- (f4.west);

    \node[desc, below=0.4cm of f1]
      {Procesar datos GTFS de Metro, Metrobús y RTP con \textit{Apimetro} para
       construir el grafo matemático de la red.};
    \node[desc, below=0.4cm of f2]
      {Medir la eficiencia actual: cuellos de botella y estaciones más vulnerables
       con \textit{VFTModel}.};
    \node[desc, below=0.4cm of f3]
      {Incorporar el corredor anillar al modelo y comparar tiempos de viaje y
       vulnerabilidad antes/después.};
    \node[desc, below=0.4cm of f4]
      {Traducir los resultados en lineamientos de integración modal, tarifaria y
       de política pública.};
  \end{tikzpicture}
  \end{center}
\end{frame}
